%% Chapter 2 — Theoretical Background
\section{Theoretical Background}
\label{sec:background}

\subsection{Superconducting Qubits}

Superconducting qubits are anharmonic LC oscillators whose nonlinearity is
provided by one or more Josephson junctions~\cite{Krantz2019}.  The Josephson
junction acts as a nonlinear inductor with energy $E_J\cos\hat{\varphi}$, where
$\hat{\varphi}$ is the superconducting phase across the junction and $E_J$ is the
Josephson energy.  Combined with a shunt capacitance $C$ (charging energy
$E_C = e^2/2C$), the system forms an artificial atom with a discrete, anharmonic
spectrum that can be addressed at the transition frequency $\omega_{01}$. Several qubit designs have been developed that exploit this anharmonicity in different parameter regimes; this thesis is primarily concerned with two: the transmon and the fluxonium.

\paragraph{Transmon.}
The transmon~\cite{Koch2007} operates deep in the regime $E_J \gg E_C$, making
its energy levels insensitive to charge noise.  Its transition frequency lies
typically in the \SIrange{4}{8}{\giga\hertz} range, and its anharmonicity
$\alpha = \omega_{12} - \omega_{01}$ is of order $-E_C/\hbar$, typically a few
hundred \si{\mega\hertz}.

\paragraph{Fluxonium.}
The fluxonium~\cite{Manucharyan2009} replaces the transmon's large capacitor with a
superinductance $L$ (inductive energy $E_L = (\Phi_0/2\pi)^2/L$, where $\Phi_0$ is
the flux quantum), threading the loop with an external flux $\Phi_z$.  The resulting
double-well potential can be tuned between a plasmon regime (near $\Phi_z = 0$) and
a fluxon regime (near $\Phi_z = \Phi_0/2$).  Fluxonium qubits offer long coherence
times and are of particular interest because their readout landscape is richer than
the transmon's: the dispersive shift $\chi$ varies strongly with flux bias, and near
the half-flux sweet spot the two eigenstates carry opposite persistent currents,
giving rise to a large state-dependent dispersive shift that can be exploited for
high-contrast readout~\cite{Manucharyan2009}.

\subsection{Circuit Quantum Electrodynamics}

In the circuit QED (cQED) architecture~\cite{Blais2004}, a qubit is coupled
dispersively to a microwave resonator.  The resonator acts both as a readout bus
and as a filter that protects the qubit from radiative decay into the measurement
chain.  The full Hamiltonian (setting $\hbar=1$) is
\begin{equation}
  \hat{H} = \hat{H}_\text{qubit}
  + \omega_r \hat{a}^\dagger \hat{a}
  + g\,\hat{n}_q (\hat{a} + \hat{a}^\dagger),
  \label{eq:jaynes_cummings}
\end{equation}
where $\omega_r$ is the resonator frequency, $\hat{a}$ ($\hat{a}^\dagger$) is the
resonator annihilation (creation) operator, $g$ is the qubit--resonator coupling
strength, and $\hat{n}_q$ is the relevant qubit operator (charge for capacitive
coupling).

\subsection{Dispersive Readout}
\label{subsec:dispersive}

When $|g_{ij}| \ll |\omega_{ij} - \omega_r|$ for all relevant qubit transitions
$i \leftrightarrow j$, second-order perturbation theory in $g$ yields a
qubit-state-dependent shift of the resonator frequency~\cite{Blais2021}:
\begin{equation}
  \hat{H}_\text{disp} = \left(\omega_r + \sum_i \chi_i \ket{i}\bra{i}\right)
  \hat{a}^\dagger \hat{a},
  \quad
  \chi_i = \sum_{j \neq i}
  \frac{|g_{ij}|^2}{\omega_{ij} - \omega_r}
  - \frac{|g_{ji}|^2}{\omega_{ji} - \omega_r},
  \label{eq:dispersive_shift}
\end{equation}
where $g_{ij} = \braket{i|g\,\hat{n}_q|j}$ and $\omega_{ij} = \omega_j -
  \omega_i$.  The textbook two-level dispersive shift $\chi = \chi_1 - \chi_0$
follows from truncating to the computational subspace, but for fluxonium-class
qubits the straddling regime (where $|g_{12}| \gg |g_{01}|$) makes this
truncation unsafe; higher levels must be included~\cite{Koch2007}.

In practice, one probes the resonator at its bare frequency $\omega_r$ and
observes a state-dependent phase shift of the transmitted or reflected signal.
For a single-qubit system, the resonator sits at $\omega_r + \chi_0$ when the
qubit is in $\ket{0}$ and at $\omega_r + \chi_1$ when in $\ket{1}$, resulting
in an IQ-plane separation that grows with the dimensionless cavity pull
$2|\chi|/\kappa$, where $\kappa$ is the resonator linewidth~\cite{Blais2004}.

\paragraph{Fluxonium dispersive readout.}
For a fluxonium qubit, the dispersive shift landscape is richer than for the
transmon because the qubit spectrum varies strongly with the external flux $\Phi_z$.
Near the half-flux point $\Phi_z = \Phi_0/2$, the double-well potential becomes
symmetric and the two lowest eigenstates carry opposite persistent currents.
The resulting state-dependent coupling to the resonator produces a dispersive shift
$\chi = \chi_1 - \chi_0$ whose magnitude depends strongly on the flux bias point,
making the readout contrast highly tunable.  This flux dependence is a key
motivation for the fluxonium readout optimisation work being pursued in related
projects within the Qilimanjaro team.

\subsection{Direct Digital Synthesis}
\label{subsec:dds}

Direct Digital Synthesis (DDS) generates analogue waveforms by accumulating a
phase register at a fixed clock rate and mapping the instantaneous phase value to
a digital amplitude via a look-up table.  A DDS numerically controlled oscillator
(NCO) produces a tone at frequency
\begin{equation}
  f_\text{out} = \frac{f_\text{clk} \cdot \Delta\phi}{2^N},
\end{equation}
where $f_\text{clk}$ is the system clock, $\Delta\phi$ is the phase increment per
clock cycle, and $N$ is the bit width of the phase accumulator.  DDS avoids the
spur and image problems associated with analogue IQ mixers because all
imperfections are digital and predictable.  In \QICK{}, the tProcessor drives DDS
engines on both the DAC (transmit) and ADC (receive) sides, enabling phase-locked
signal generation and coherent demodulation~\cite{Stefanazzi2022}.

\subsection{State Discrimination}

Following dispersive readout, the measured IQ quadrature amplitudes $(I, Q)$ form
two clusters in the complex plane — one for $\ket{0}$ and one for $\ket{1}$.
Single-shot state discrimination assigns each measurement to one of these
clusters.  The simplest approach is a linear threshold: a hyperplane in the IQ
plane that minimises the total misclassification probability, which for Gaussian
clusters of equal covariance is the perpendicular bisector of the line connecting
the centroids~\cite{Gambetta2007}.  The single-shot readout fidelity is
\begin{equation}
  \mathcal{F} = 1 - \frac{1}{2}\left(p(1|0) + p(0|1)\right),
\end{equation}
where $p(1|0)$ is the probability of assigning $\ket{1}$ given that the qubit was
prepared in $\ket{0}$, and vice versa.  More sophisticated classifiers — quadratic
discriminant analysis, neural networks — can improve fidelity when the clusters are
non-Gaussian or when multi-level leakage is present.
